\section{Truncated Bases: Captured Energy Governs}\label{sec:truncation}

\paragraph{Why the wavelet premise was misconceived.}
Which property of compactly supported orthogonal dyadic wavelets does a frozen
synthesis basis exploit? None. The classical superiority of wavelets is a theorem
about \emph{nonlinear} best-$m$-term approximation, in which the retained
coefficients are chosen per signal. Here the rows of $\Bf$ are fixed at build time
and shared across every series in the dataset. That is \emph{linear}
approximation, and under linear approximation the optimal fixed subspace is the
Karhunen--Lo\"eve basis, which the DCT approximates for stationary covariance.
Localisation---the property wavelets are chosen for---becomes a liability under
fixed truncation, because each retained translate covers only part of the time
axis. The time--frequency versus time--scale axis is therefore the wrong axis for
this architecture. What governs the outcome is how much of the target's energy
lies in the span of the retained rows, whatever produced those rows.

\paragraph{The band-boundary case.}
M4-Monthly makes the mechanism concrete. Its dyadic bands sit at nine rows. A
band-ordered wavelet truncated to six rows captures $0.658$ of the target energy
and scores $73.380$ \smape; at nine rows it captures $0.996$ and scores $13.850$,
level with a rank-matched DCT at $13.720$. Three basis rows are worth $59.5$
\smape. Nothing about the wavelet changed between the two cells; the truncation
either respects the band boundary or it does not.

\paragraph{The one-third-rank ordering.}
Table~\ref{tab:hourly_third_rank} orders four bases on M4-Hourly at one-third
forecast rank ($\df = 16$ of $\Ho = 48$), over ten seeds at identical parameter
counts. The Karhunen--Lo\"eve basis that linear approximation theory names as
optimal leads at $10.054$, the DCT follows at $10.287$, and a band-ordered
Daubechies~3 basis scores $10.587$. The same wavelet under SVD row ordering
collapses to $73.946$. The ordering follows captured energy exactly, and the gap
between the two wavelet rows is an ordering effect, not a family effect.

\begin{table}[t]
  \centering
  \caption{M4-Hourly at one-third forecast rank ($\df = 16$ of $\Ho = 48$), ten
  seeds, identical parameter counts. The trained ordering follows captured
  energy exactly, including where capture differs only in the fourth decimal.}
  \label{tab:hourly_third_rank}
  \begin{tabular}{lcc}
    \toprule
    Basis & Captured energy & \smape \\
    \midrule
    Karhunen--Lo\"eve & $0.999981$ & $10.054$ \\
    DCT & $0.999964$ & $10.287$ \\
    Daubechies 3, band-ordered & $0.999885$ & $10.587$ \\
    Daubechies 3, SVD-ordered & $0.347765$ & $73.946$ \\
    \bottomrule
  \end{tabular}
\end{table}

\paragraph{The capture law.}
Traffic-96 replicates the ordering across five bases at identical parameter
counts (Table~\ref{tab:traffic_capture}). Captured energy falls from $0.9714$ to
$0.3035$ while \smape{} rises from $29.78$ to $120.98$, without a single
inversion (Spearman $-1.0$, exact $p = 0.017$, the floor at $n = 5$; run-level
$\rho = -0.98$ over 25 runs, $p = 9.0\mathrm{e}{-18}$). A quantity computable
from the basis and the data alone, before any training, predicts the trained
ordering exactly.

\begin{table}[t]
  \centering
  \caption{Traffic-96, five bases at identical parameter counts, ordered by
  captured energy. Endpoint values are shown; the three intermediate rows are
  marked with em-dashes, and the full per-basis and per-run values ship with the
  reproducibility bundle. The ordering has no inversion (Spearman $-1.0$, exact
  $p = 0.017$, the floor at $n = 5$; run-level $\rho = -0.98$ over 25 runs,
  $p = 9.0\mathrm{e}{-18}$).}
  \label{tab:traffic_capture}
  \begin{tabular}{cc}
    \toprule
    Captured energy & \smape \\
    \midrule
    $0.9714$ & $29.78$ \\
    --- & --- \\
    --- & --- \\
    --- & --- \\
    $0.3035$ & $120.98$ \\
    \bottomrule
  \end{tabular}
\end{table}

\paragraph{A withdrawn claim.}
We withdraw an earlier claim, made after submission, that a wavelet basis
finishes behind random draws. That claim measured our own orthogonalisation: the
construction we shipped orthogonalised the raw wavelet matrix by SVD and ordered
the rows by singular value, which mixes every wavelet function together and
destroys the dyadic coarse-to-fine ordering, so truncation retained an arbitrary
subspace rather than the coarse scales. Rebuilt with QR---triangular, so the
dyadic ordering survives---the same wavelet lands within $0.30$ \smape{} of a
DCT. The failure was in the row ordering, not in the wavelet.

\paragraph{Tiering, resolved.}
What governs tiering is spatial coverage: a localised basis covers the time axis
in pieces, so cutting off inside one of its scale bands leaves part of the
horizon unreachable---the Monthly band-boundary case above. What does not
survive is the claim that tiering helps. A selection fixed at build time cannot
adapt per series, and a direct test agrees: on M4-Hourly, plain truncation by
rank gives $8.788$ against $8.926$ for offset tiering with a Fourier basis, and
$8.863$ against $8.885$ with a wavelet. The offset mechanism adds nothing in
either family.

Captured energy says how accurate a truncated basis will be. Whether the model
will train at all is a separate axis, governed by a different quantity;
Section~\ref{sec:trainability} takes it up.
